\section{Projection publique du registre interne des routes}

Chaque route possède \(53\) champs dans le registre unifié et \(57\) champs dans le
registre enrichi. La fiche scientifique publie \(49\) champs. Cette
projection n'élimine pas d'information : les états de cohérence de la route,
du raccordement et du prolongement par tours, ainsi que les \(5\) composantes
brutes de la signature, demeurent dans les registres complets et sont réunis dans les
invariants publiés. Par conséquent, l'application de publication est une
projection typée
\[
  \mathcal R_{57}\longrightarrow\mathcal F_{49},
\]
accompagnée du registre intégral \(\mathcal R_{53}\) et du registre
enrichi \(\mathcal R_{57}\) ; ce n'est pas une réduction du domaine scientifique.

\section{Correspondance entre le domaine généalogique et les inventaires métrologiques}

Le domaine matériel est construit au moyen de \(81\) cellules, de \(56\) familles de
routes, de \(13\) multisections et de \(324\) routes orientées. Son déploiement intégral
se trouve dans
\hyperref[chap:atlas-genealogico-exhaustivo-324-rutas]{l'atlas de \(324\)
routes}, avec une fiche complète et \(49\) champs par route. Les recensements externes
de confrontation sont conservés dans
\hyperref[chap:inventario-universal-471-masas-384-anchuras]{l'inventaire de
\(471\) masses et de \(384\) largeurs}. L'atlas construit les fibres et leurs
opérateurs ; les recensements ne sélectionnent pas rétrospectivement une route ni ne
redéfinissent le domaine de la loi.
